\documentclass[12pt]{book}
\usepackage{amsmath}
\usepackage{amssymb}
\usepackage{geometry}
\geometry{margin=1in}

\title{Volume 02: Field Theory, Mechanics, and Mathematical Structure \\ Book 01: Continuum Mechanics from Matrix Flow}
\author{James C. Harvey}
\date{December 2025}

\begin{document}
\maketitle
\tableofcontents

\chapter{Continuum Mechanics from Matrix Flow}

\section*{Abstract}
This book derives continuum mechanics from SDT by treating fluids as matrix-flow phenomena that transport
matter through flux coupling. Navier--Stokes emerges from momentum transfer through boundary locking,
viscosity is a geometric consequence of matrix contact friction, and turbulence is deterministic chaos
in flux topology. Every equation is written in SDT terms, with explicit constants and dimensional checks.

\section{Foundational Framing}
Continuum mechanics in SDT does not begin with mass and viscosity as primitives. It begins with space,
matter, movement, and now. Space is the matrix, flux is its directional propagation, and the medium is
the collective behavior of matrix units (spations). Fluid motion is a macroscopic expression of flux
transport coupled to matter boundaries.

\section{Matrix Flux and Momentum Transfer}
Define the matrix momentum flux:
\begin{equation}
\mathbf{j}_m = \rho_m \mathbf{v}_m,
\end{equation}
where $\rho_m$ is the matrix density and $\mathbf{v}_m$ is the matrix flow velocity. Matter boundaries
exchange momentum with the matrix through locking contacts. The momentum transfer rate is
\begin{equation}
\frac{d\mathbf{p}}{dt} = \lambda \mathbf{j}_m \cdot \mathbf{n}\, A,
\end{equation}
where $\lambda$ is the locking efficiency, $\mathbf{n}$ is the surface normal, and $A$ is the contact
area.

\section{Pressure Field from CMB Flux}
The pressure field in the matrix is sourced by the CMB flux:
\begin{equation}
P(\mathbf{r}) = P_{\text{CMB}} + \Delta P(\mathbf{r}),
\end{equation}
with $P_{\text{CMB}} = 2.036 \times 10^{-2}\ \text{Pa}$ at atomic/molecular scale. The pressure gradient
is the driver of macroscopic flow:
\begin{equation}
\mathbf{f}_{\text{pressure}} = -\nabla P = -\nabla(\Delta P).
\end{equation}

\section{Navier--Stokes from Matrix Mechanics}
The standard Navier--Stokes form emerges from SDT momentum transfer:
\begin{equation}
\rho \frac{D\mathbf{v}}{Dt} = -\nabla P + \mu \nabla^2 \mathbf{v} + \mathbf{f},
\end{equation}
where $\rho$ is the matter density, $\mathbf{v}$ is the matrix flow velocity, and $\mathbf{f}$ represents
body forces from pressure gradients. The dynamic viscosity $\mu$ is not a primitive; it arises from
matrix contact friction:
\begin{equation}
\mu = \lambda \frac{K_{\text{bulk}} \ell_P^2}{c}.
\end{equation}
Dimensional check:
\begin{equation}
[\mu] = [K_{\text{bulk}} \ell_P^2 / c] = \text{Pa}\cdot \text{s} = \text{kg/(m·s)}.
\end{equation}

\section{Numerical Form of Viscosity}
With canonical constants,
\begin{equation}
\mu = \lambda \frac{(4.6 \times 10^{113}) (1.616255 \times 10^{-35})^2}{2.99792458 \times 10^8},
\end{equation}
which makes the viscosity an explicit numeric function of locking efficiency $\lambda$.

\section{Incompressibility Constraint}
The matrix is incompressible at the macroscopic scale of flow:
\begin{equation}
\nabla \cdot \mathbf{v} = 0.
\end{equation}
This condition is a geometric constraint on flux distribution, not a separate physical law.

\chapter{Vorticity and Flux Rotation}

\section*{Abstract}
Vorticity is the rotational content of matrix flow. It is the local measure of circulation density and
sets the topology of eddies, wakes, and shear layers.

\section{Vorticity Definition}
Define vorticity as
\begin{equation}
\boldsymbol{\omega} = \nabla \times \mathbf{v}.
\end{equation}
When $\boldsymbol{\omega} \neq 0$, the matrix flow contains rotational structure, which is stabilized
by occlusion geometry and boundary proximity.

\section{Vorticity Transport}
Taking the curl of the flow equation yields the vorticity transport form:
\begin{equation}
\frac{D\boldsymbol{\omega}}{Dt} = (\boldsymbol{\omega}\cdot\nabla)\mathbf{v} + \nu \nabla^2 \boldsymbol{\omega},
\end{equation}
where $\nu = \mu/\rho$ is the kinematic viscosity. In SDT, $\nu$ inherits geometric meaning through
$\mu = \lambda K_{\text{bulk}}\ell_P^2/c$.

\chapter{Dimensionless Structure and Flow Regimes}

\section*{Abstract}
Regime classification is a geometric ratio of inertial flux to locking dissipation. This chapter derives
the SDT analogs of Reynolds and Mach numbers.

\section{Flux Reynolds Number}
Define the Reynolds-like ratio:
\begin{equation}
\mathrm{Re}_m = \frac{\rho v L}{\mu} = \frac{\rho v L c}{\lambda K_{\text{bulk}}\ell_P^2}.
\end{equation}
Large $\mathrm{Re}_m$ indicates inertial dominance and turbulence; small $\mathrm{Re}_m$ indicates
locking-dominated laminar flow.

\section{Regime Threshold Estimation}
For a characteristic scale $L$ and velocity $v$, the transition condition
\begin{equation}
\mathrm{Re}_m \approx 10^3
\end{equation}
provides a geometric boundary between laminar and turbulent flux regimes.

\section{Flux Mach Number}
The compressibility ratio is
\begin{equation}
\mathrm{Ma}_m = \frac{v}{c}.
\end{equation}
Because the matrix supports propagation at $c$, $\mathrm{Ma}_m$ measures the proximity to compressible
behavior.

\chapter{Boundary Layer Scaling}

\section*{Abstract}
Boundary layers are the spatial record of locking gradients. This chapter derives their thickness from
balance between inertial transport and matrix contact dissipation.

\section{Thickness Estimate}
Let $U$ be free-stream speed and $x$ the downstream coordinate. The characteristic thickness is
\begin{equation}
\delta(x) \sim \sqrt{\frac{\nu x}{U}}.
\end{equation}
In SDT, $\nu$ is geometric; therefore $\delta$ is a geometric consequence of matrix locking length scales.

\chapter{Turbulence as Deterministic Chaos}

\section*{Abstract}
Turbulence is deterministic chaos in matrix flux. It emerges from nonlinear coupling in the flow equation
and is therefore not random.

\section{Nonlinear Flux Coupling}
The convective term $(\mathbf{v}\cdot\nabla)\mathbf{v}$ produces sensitive dependence on initial conditions.
The turbulent cascade is a hierarchy of flux structures arising from geometric instability of flow paths.

\section{Eddies as Flux Topology}
Eddies are not separate entities; they are closed flux loops stabilized by local occlusion geometry and
boundary proximity.

\section{Energy Cascade}
The kinetic energy per unit mass is $k = \frac{1}{2}\langle v^2\rangle$. The cascade rate $\varepsilon$ is
the flux of kinetic energy through scales:
\begin{equation}
\varepsilon \sim \frac{u_\ell^3}{\ell},
\end{equation}
where $u_\ell$ is characteristic velocity at scale $\ell$. In SDT, this cascade is a deterministic
redistribution of matrix flux guided by occlusion topology.

\chapter{Boundary Layers from Locking Gradients}

\section*{Abstract}
Boundary layers arise where locking efficiency varies with distance from matter boundaries. Velocity
gradients are therefore a direct map of $\lambda$ gradients in the matrix.

\section{Locking Efficiency Gradient}
\begin{equation}
\nabla \lambda \neq 0 \quad \Rightarrow \quad \nabla \mathbf{v} \neq 0.
\end{equation}
The thickness of the boundary layer is the scale over which locking transitions from high to low.

\chapter{Compressible Flow from Density Variations}

\section*{Abstract}
Compressibility is permitted when local matrix density varies with pressure. These variations are not
fundamental compressions of the matrix, but local redistribution of flux density.

\section{Density Variation}
\begin{equation}
\rho_m = \rho_m(P),
\end{equation}
so that changes in $P$ induce local variations in effective density and allow compressible behavior in
macroscopic flow.

\section{Flux Equation of State}
For matrix-supported flow, pressure and density are linked through stiffness:
\begin{equation}
\Delta P = K_{\text{bulk}} \frac{\Delta \rho_m}{\rho_m}.
\end{equation}
This relation governs compressible wave behavior and shock formation.

\chapter{Plasma as Charged Vortex Interaction}

\section*{Abstract}
Plasma is a regime of charged vortex interaction where flux geometry organizes large-scale collective
behavior. This chapter positions plasma physics as a special case of matrix flow with strong coupling to
boundary circulation.

\section{Charged Vortex Coupling}
The coupling between charged vortices modifies the effective pressure gradient:
\begin{equation}
\nabla P_{\text{eff}} = \nabla P + \nabla P_{\text{wake}},
\end{equation}
where $P_{\text{wake}}$ is the helical wake contribution from circulation.

\chapter{The Narrative of Continuum Mechanics}

\section*{Abstract}
This chapter ties the derivation into a single narrative: flux produces pressure; pressure produces
gradients; gradients produce movement; movement produces structure.

\section{Flux $\rightarrow$ Pressure}
\begin{equation}
\Pi(\mathbf{r}) = \int_{\Omega} I_{\text{CMB}}(\hat{\mathbf{n}})\left[1 - E(\mathbf{r},\hat{\mathbf{n}})\right]\, d\Omega.
\end{equation}

\section{Pressure $\rightarrow$ Movement}
\begin{equation}
\mathbf{F} = -\nabla \Pi(\mathbf{r}).
\end{equation}

\section{Movement $\rightarrow$ Flow Law}
The Navier--Stokes equation is the macroscopic form of this chain.

\end{document}
